\section{Conclusion}

This paper presented ASID, an adaptive sparse inference framework for online multimodal fault diagnosis in cloud microservice systems. ASID addresses the fundamental tension between diagnostic accuracy and inference efficiency by integrating three key designs: modality-specific encoding with deviation-aware temporal modeling to construct anomaly-sensitive representations from metrics, logs, and traces; \lb{a frozen truncated GPT-2 backbone with parameter-efficient adaptation} for sequential modeling; and \lb{a service-aware sparse expert routing mechanism with a confidence-adaptive budget that activates one or two lightweight adapters according to runtime service states}.

\lb{Experiments on the Eadro-SN benchmark demonstrate that ASID achieves an F1 score of 0.9838 with zero deadline misses under a 25\,ms tight deadline, a p99 latency of 16.69\,ms, and a GPU peak memory of 149.43\,MB under the optimized serving path. The confidence-adaptive budget avoids unnecessary expert evaluations on confident windows while preserving diagnostic quality. Compared with the Eadro baseline, ASID improves F1 by 5.02 percentage points. Compared with XGBoost, ASID improves F1 by 5.76 percentage points while reducing p99 latency from 40.15\,ms to 16.69\,ms under full online replay.} Ablation studies confirm that the multimodal representation, sparse expert adapter, service-conditioned routing, and guarded online decision strategy contribute to the final performance. \lb{Supplementary anomaly-detection results on MSDS and RE2-TT show that ASID also transfers beyond the primary benchmark. Timing evaluations further show zero miss rate under deadline budgets of 25\,ms and above, zero miss rate at up to 6.7 times the nominal arrival load in trace-driven queue analysis, and low GPU memory consumption in the optimized serving path.} A supplementary RCAEval RE2-TT study further shows that the same deviation representation can support service-level root-cause ranking through an optional root-victim branch.

\lb{These results demonstrate that adaptive sparse inference is a practical and effective paradigm for deploying intelligent fault diagnosis in latency-sensitive cloud environments under online replay. Future work will focus on concurrent deployment tests beyond trace-driven queue analysis and more elastic service-prior handling for newly introduced or renamed services.}
